\begin{abstract}
Malicious news sources can change their presentation without changing the probability that their publications are false. This study evaluates such strategies in X through a survey-informed agent-based model. Five hundred heterogeneous users repeatedly select and repost publications from four source types according to accumulated evaluations, visibility, social recommendations, and memory. Malicious sources imitate engagement, proximity, information-quality, and credibility attributes of traditional media. We conduct 2,010 seeded runs across four experiments concerning strategy type, campaign timing and duration, reach restrictions, and behavioural sensitivity. Strategies increased the malicious source's audience, especially when all attribute groups were imitated, but generally produced small changes in the overall share of false reposts under truth-sensitive recommendations. Reach restrictions reduced malicious-source participation and eliminated strategy effects at zero direct reach. Truth-sensitive recommendations attenuated the additional false-repost effect observed when recommendations were disabled. Sensitivity analysis showed that audience-capture conclusions were robust, whereas false-repost effects depended on memory persistence. The findings distinguish source success from societal misinformation exposure and suggest evaluating reach governance and recommendation mechanisms jointly. The model is a policy-oriented computational laboratory rather than a predictive representation of X.
\end{abstract}

\noindent\textbf{Keywords:} agent-based modelling; fake news; misinformation; news sources; social media; scenario analysis
